\chapter{Experimental Protocol and Metrics}
\label{ch:exp_protocol_metrics}

This chapter defines the evaluation protocol shared by every experiment in
this thesis.  Section~\ref{sec:pipeline_orchestration} outlines the
eight-stage pipeline from configuration to evaluation.
Section~\ref{sec:splits} defines the train/test split.
Section~\ref{sec:lag_selection} motivates the MVAR lag order via an
information-theoretic study.
Section~\ref{sec:metrics} specifies all evaluation metrics and the forecast
horizon.

%% -------------------------------------------------------------------
\section{Pipeline orchestration}
\label{sec:pipeline_orchestration}

Every experiment is driven by a single YAML configuration file (one per
regime) and executed end-to-end by the pipeline described in
Figure~\ref{fig:pipeline_stages}.
The pipeline proceeds through the eight sequential stages shown in
Figure~\ref{fig:pipeline_stages}.

\begin{figure}[htbp]
\centering
\begin{tikzpicture}[
    node distance=0.35cm and 0.0cm,
    stage/.style={
      draw, rounded corners=3pt, minimum width=13cm,
      minimum height=0.65cm, font=\footnotesize, align=center,
      fill=blue!6
    },
    arr/.style={-{Stealth[length=3pt]}, semithick}
  ]
  \node[stage]                  (s1) {\textbf{1.}~Generate training data --- 340 Vicsek--Morse trajectories, 4~IC families};
  \node[stage, below=of s1]     (s2) {\textbf{2.}~Build shared POD basis --- shift-align, $\sqrt{\cdot}$-transform, truncated SVD $\to$ $r$ modes};
  \node[stage, below=of s2]     (s3) {\textbf{3.}~Build latent dataset --- project $\to$ $y(t)\in\mathbb{R}^r$, subsample$\,{=}\,3$, construct $(X,Y)$ pairs};
  \node[stage, below=of s3]     (s4) {\textbf{4.}~Train MVAR \& LSTM on shared latent dataset (Chapter~\ref{ch:latent_rom})};
  \node[stage, below=of s4]     (s5) {\textbf{5.}~Generate test data --- 4 held-out trajectories (one per IC family)};
  \node[stage, below=of s5]     (s6) {\textbf{6.}~Evaluate MVAR \& LSTM --- condition, forecast, lift, compute metrics};
  \node[stage, below=of s6]  (s7) {\textbf{7.}~WSINDy PDE discovery --- same train split, raw-field weak systems, MSTLS};
  \node[stage, below=of s7]  (s8) {\textbf{8.}~WSINDy evaluation --- identification diagnostics, cross-scale stability selection, targeted closure stress test};
  %% arrows
  \draw[arr] (s1) -- (s2);
  \draw[arr] (s2) -- (s3);
  \draw[arr] (s3) -- (s4);
  \draw[arr] (s4) -- (s5);
  \draw[arr] (s5) -- (s6);
  \draw[arr] (s6) -- (s7);
  \draw[arr] (s7) -- (s8);
\end{tikzpicture}
\caption{Eight-stage pipeline.
  Stages~1--6 run the EF-ROM branch~\cite{alvarez2025eqfree};
  stages~7--8 run WSINDy~\cite{messenger2021wsindy} PDE discovery and evaluation.}
\label{fig:pipeline_stages}
\end{figure}

\noindent
Stages~1 and~5 share a common upstream: the same
Vicsek--Morse~\cite{vicsek1995vicsek,dorsogna2006morse} simulations,
KDE-based coarse-grained fields, and train/test split.
Stages~2--6 build a shift-aligned~\cite{reiss2018shifted_pod} POD latent
dataset and learn MVAR and LSTM~\cite{hochreiter1997lstm} evolution maps;
stages~7--8 bypass POD and apply
MSTLS~\cite{minor2025mstls}/stability-selection~\cite{meinshausen2010stability}
weak-form regression on the same trajectories, with
ETDRK4~\cite{kassam2005etdrk4} rollout performed only as a targeted
closure diagnostic on selected identified PDEs (not as the primary
evaluation protocol).  All artefacts and metrics are written to a
self-contained experiment directory.

%% -------------------------------------------------------------------
\section{Train/test splits}
\label{sec:splits}

Training and test data are drawn from the same parameter regime but use
disjoint initial conditions.

\noindent\textbf{Training set.}
Each regime generates 340 runs from four IC families
(Table~\ref{tab:ic_split}, Chapter~\ref{ch:experimentation}).
For the EF-ROM branch, all runs contribute to the shared POD basis and latent
dataset.  For WSINDy, the same training trajectories are used in physical
space: density and auxiliary fields are computed on each run and then
converted into weak-form regression constraints.

\noindent\textbf{Test set.}
Four runs---one per IC family---are generated at held-out IC parameters
(e.g.\ a Gaussian centred at $(12.5,12.5)$ rather than the training grid
positions) and simulated for $T_{\mathrm{test}}$ (Table~\ref{tab:fixed_params}).
They are never seen during POD construction, latent-model fitting, or WSINDy
regression.  The same held-out trajectories define the test set for both
branches.

\noindent\textbf{Forecasting protocol.}
All models are evaluated in closed-loop (iterated single-step) mode.
The effective forecast start is
$t_{\mathrm{fc}} = \max(t_{\mathrm{fc}}^{\mathrm{req}},\, w_{\max}\Delta t_{\mathrm{rom}})$,
where $w_{\max}$ is the largest enabled latent lag.
In baseline configs $w_{\mathrm{MVAR}}=w_{\mathrm{LSTM}}=5$, giving
$t_{\mathrm{fc}}=0.6$\,s.
WSINDy evaluation uses the same $t_{\mathrm{fc}}$ but initialises from
physical fields rather than a lagged latent window.

%% -------------------------------------------------------------------
\section{Lag selection}
\label{sec:lag_selection}

The MVAR lag order $w$ determines how many past latent frames the model
uses for prediction.  We evaluate BIC, HQIC, and AIC~\cite{wasserman2006all} on unrestricted
VAR($w$) models for $w=1,\dots,100$ on every regime.
For reference, the information-theoretically optimal lags on a
representative DO catalogue regime (DO\_CS01, collective swarm; see
Appendix~\ref{app:full_regime_catalogue}) range from
$w^*_{\mathrm{BIC}}=23$ to $w^*_{\mathrm{AIC}}=50$.
Despite these large information-criterion optima, we fix
$w_{\mathrm{MVAR}}=5$ for the baseline systematic comparison.
This is a deliberate design trade-off between three considerations:
\begin{enumerate}[nosep]
\item \textbf{Cross-regime comparability.}  A single lag order applied
      uniformly to all nine benchmark regimes ensures that performance
      differences reflect the dynamics rather than per-regime
      hyperparameter tuning.
\item \textbf{Computational cost.}  The parameter count of an MVAR($w$)
      model scales as $w r^2$; at $r=19$, increasing $w$ from~5 to~23
      would multiply the number of coefficients by~4.6$\times$.
\item \textbf{Diminishing returns after preprocessing.}  After
      shift-alignment and the $\sqrt{\cdot}$-transform in our pipeline,
      the latent trajectories become approximately linear, so additional
      lags contribute less new information than in the raw
      (unconditioned) setting where the BIC/AIC criteria were computed.
\end{enumerate}
No formal lag sweep over the deployed ridge-regularised MVAR is presented
in this thesis; validating whether $w{>}5$ yields measurable forecast
gains is left to future work (Section~\ref{sec:future_work}).
We note that the BIC analysis was performed on a single DO catalogue
regime; the information-criterion optima may differ for the NDYN
benchmark regimes, so $w{=}5$ is a conservative lower bound rather
than a per-regime optimum.
The choice $w{=}5$ should therefore be interpreted as a practical,
not individually optimal, setting.
The LSTM lag is treated as a separate hyperparameter; whenever
$w_{\mathrm{LSTM}}\neq w_{\mathrm{MVAR}}$, the common forecast
start is shifted according to the largest enabled lag.

%% -------------------------------------------------------------------
\section{Evaluation metrics}
\label{sec:metrics}

All metrics compare predicted densities $\hat{\boldsymbol{\rho}}^n$ against
ground truth $\boldsymbol{\rho}^n$ over the forecast window
$\mathcal{T}_{\mathrm{fc}} = \{n_0, n_0{+}1, \dots, n_1\}$, where $n_0$
is the forecast start index and $n_1$ is the terminal index.
Let $d = N_x N_y$ denote the number of spatial grid points.

\noindent\textbf{Coefficient of determination ($R^2$).}
\label{subsec:r2_time}
The per-snapshot $R^2$ is
\begin{equation}
  R^2(n) \;=\; 1 \;-\;
  \frac{\sum_{j=1}^{d}\bigl(\rho_j^n - \hat\rho_j^n\bigr)^2}
       {\sum_{j=1}^{d}\bigl(\rho_j^n - \bar\rho^n\bigr)^2},
  \qquad
  \bar\rho^n = \frac{1}{d}\sum_{j=1}^d \rho_j^n,
  \label{eq:r2_snapshot}
\end{equation}
where the denominator uses the spatial mean of the ground-truth snapshot at
time~$n$.  We report the mean over the forecast window,
$R^2_{\mathrm{fc}} = |\mathcal{T}_{\mathrm{fc}}|^{-1}\sum_{n\in\mathcal{T}_{\mathrm{fc}}} R^2(n)$,
as the primary scalar metric.
Because the null model in~\eqref{eq:r2_snapshot} is the per-snapshot
spatial mean~$\bar\rho^n$---a single scalar per frame---$R^2$ values
near~$1$ indicate that the predictor tracks the full spatial structure,
not merely the global density level.  For nearly uniform density fields,
however, the denominator becomes small and $R^2$ can be dominated by
numerical noise; the relative-error norms below provide a complementary
check in such regimes.

\noindent\textbf{Relative $\ell^p$ errors.}
\label{subsec:rel_errors}
Three relative error norms are
averaged over the
forecast window:
\begin{align}
  \mathrm{RelErr}_1(n) &= \frac{\|\boldsymbol{e}^n\|_1}{\|\boldsymbol{\rho}^n\|_1 + \epsilon},
  &
  \mathrm{RelErr}_2(n) &= \frac{\|\boldsymbol{e}^n\|_2}{\|\boldsymbol{\rho}^n\|_2 + \epsilon},
  &
  \mathrm{RelErr}_\infty(n) &= \frac{\|\boldsymbol{e}^n\|_\infty}{\|\boldsymbol{\rho}^n\|_\infty + \epsilon},
  \label{eq:rel_errors}
\end{align}
where $\boldsymbol{e}^n = \boldsymbol{\rho}^n - \hat{\boldsymbol{\rho}}^n$
and $\epsilon = 10^{-12}$ guards against division by zero.

\noindent\textbf{Mass conservation.}
\label{subsec:mass_metric}
The discrete total mass at time $n$ is
$m^n = (\Delta x\,\Delta y)\,\mathbf{1}^\top\boldsymbol{\rho}^n$ (and
analogously $\hat m^n$ for the prediction).  We track $m^n$ and $\hat m^n$
over the forecast window and report the maximum relative mass violation:
\begin{equation}
  \Delta m_{\max} \;=\;\max_{n\in\mathcal{T}_{\mathrm{fc}}}
    \frac{|\hat m^n - m^n|}{m^{n_0}}.
  \label{eq:mass_violation}
\end{equation}
The pipeline optionally applies a post-hoc mass correction step,
either the global rescaling
\eqref{eq:mass_postprocess_scale} or the simplex projection
\eqref{eq:mass_postprocess_simplex}, to match the initial mass.  In the
preprocessing ablations this had little effect on MVAR but improved LSTM
rollout robustness and reduced $\Delta m_{\max}$ in practice.

\noindent\textbf{Spatial order parameter.}
\label{subsec:spatial_order_metric}
To assess whether the model preserves the spatial structure of the density
field, we track the spatial standard deviation
$\sigma_\rho(n) = \mathrm{std}(\boldsymbol{\rho}^n)$ defined in
Eq.~\eqref{eq:spatial_order} (Chapter~\ref{ch:background}).
A high $\sigma_\rho$ indicates concentrated density; a low value
indicates a nearly uniform field.  Together with the agent-level order
parameters defined in Section~\ref{sec:order_parameters}, this diagnostic
helps distinguish regime-level structural failures that may not be fully
visible from density-space~$R^2$ alone.

\noindent\textbf{Metric hierarchy by branch.}
For the EF-ROM branch, forecast accuracy metrics---particularly
$R^2_{\mathrm{fc}}$ and the relative errors~\eqref{eq:rel_errors}---are
primary evaluation targets.  For the WSINDy branch, these same metrics are
reported for completeness and as consistency checks under ETDRK4 rollout,
but the primary criteria are identification-oriented: sparsity of the
discovered PDE, coefficient stability under cross-scale selection, agreement
with expected regime-specific mechanisms, and weak-form residual quality.
Common metrics provide a shared diagnostic vocabulary, but the decision
criteria differ between the two branches.

\noindent\textbf{Parameter counts.}
\label{subsec:param_counts}
At the baseline comparison settings
($r{=}19$, $w_{\mathrm{MVAR}}{=}w_{\mathrm{LSTM}}{=}5$, hidden width
$h{=}128$, $L{=}2$ LSTM layers), the MVAR has
$p_{\mathrm{MVAR}} = w\,r^2 + r = 1{,}824$ trainable parameters, while the
LSTM has $p_{\mathrm{LSTM}} = 211{,}091$---a roughly $116\times$
ratio.  This count compares raw learnable parameters, not effective
degrees of freedom: LSTM weight sharing across time steps and the
recurrent bottleneck mean that the effective capacity gap is smaller
than $116\times$.  Full derivations appear in
Chapter~\ref{ch:latent_rom}, \S\ref{subsec:mvar_identifiability}
and~\S\ref{sec:lstm_surrogate}.

\medskip\noindent
From this point onward, all branch-level results should be read under
this common upstream protocol but evaluated against branch-specific
decision criteria: forecast accuracy for the EF-ROM, and
identification quality for WSINDy.
